\section{Results}
\label{ch:results}

\subsection{Evidence overview}

The results examine whether the agents complete the common racing task and whether the application preserves interpretable evidence. \Cref{tab:agent-results} separates application database records from repeated frozen-policy evaluations: the former demonstrate integrated behaviour, while the latter provide stronger reliability evidence. All use the common \texttt{g-track-3}, car and one-lap scenario described in \cref{ch:testing}.

\begin{table}[htbp]
  \centering
  \caption{Agent performance under the common one-lap \texttt{g-track-3} scenario. Database records and repeated fixed-policy evaluations are intentionally reported as separate evidence sources.}
  \label{tab:agent-results}
  \small
  \begin{tabularx}{\textwidth}{p{0.22\textwidth}p{0.18\textwidth}p{0.07\textwidth}p{0.08\textwidth}p{0.11\textwidth}YY}
    \toprule
    Agent & Evidence source & Attempts & Completed & Completion rate & Best completed lap & Median completed lap \\
    \midrule
    Rule-Based Anti-Spin & Application database & 2 & 2 & 100\% & 98.394 s & 99.142 s \\
    Map-Aware Racing-Line & Application database & 15 & 15 & 100\% & 90.958 s & 92.056 s \\
    Dyna-Q Learning & Application database & 3 & 1 & 33.3\% & 305.562 s & 305.562 s \\
    TD3 Scratch & Application database & 6 & 0 & 0\% & N/A & N/A \\
    Agent 7 N-step TD3 v3 & Repeated evaluation & 20 & 18 & 90\% & 103.230 s & 104.220 s \\
    Agent 7 N-step TD3 v4 & Repeated evaluation & 10 & 10 & 100\% & 103.386 s & 103.887 s \\
    Agent 8 Sensor-Only TD3 & Repeated evaluation & 40 & 36 & 90\% & 83.038 s & 83.038 s \\
    \bottomrule
  \end{tabularx}
\end{table}

\subsection{Engineered baselines and early learning results}

\Cref{fig:completion-reliability} reports completion proportion, its 95\% Wilson interval and the raw denominator, so 2/2 is not treated as equally certain as 36/40. The Map-Aware Racing-Line Agent completed all 15 stored attempts, with no off-track events and a 92.056-second median lap. It is therefore a substantive benchmark: explicit route knowledge and engineered control were both fast and reliable on the known circuit. The rule-based agent also completed 2/2, but that small batch supports only a weak reliability claim.

\begin{figure}[htbp]
  \centering
  \includegraphics[width=0.90\textwidth]{figure_10_1_completion_reliability.png}
  \caption{Observed lap-completion proportions with 95\% Wilson intervals. Labels retain completed and attempted laps, so small perfect batches are not interpreted as equally certain as larger repeated evaluations. Source: project evaluation records.}
  \label{fig:completion-reliability}
\end{figure}

The early learning results reject the assumption that a more complex algorithm automatically wins. Dyna-Q completed 1/3 at 305.562 seconds, while scratch TD3 completed 0/6. These results expose the limits of coarse discretisation and immature continuous control, and motivate repeated completion rather than a single fast event as the final selection criterion.

\subsection{Learned-policy performance}

\Cref{fig:lap-time-distribution} retains every completed-lap observation rather than reducing each group to its fastest result; failed attempts remain in \cref{fig:completion-reliability}. Agent~7 established repeated learned completion: v3 completed 18/20 at a 104.220-second median, and v4 completed 10/10 at 103.887 seconds. Both batches were slower than the map-aware benchmark.

\begin{figure}[htbp]
  \centering
  \includegraphics[width=0.90\textwidth]{figure_10_2_lap_time_distribution.png}
  \caption{Completed-lap time distributions. Individual points are retained alongside the median and interquartile range; the plot excludes failures because those attempts have no completed-lap time. Dyna-Q is reported in \cref{tab:agent-results} but omitted here because it has only one completed lap. Source: project evaluation records.}
  \label{fig:lap-time-distribution}
\end{figure}

\begin{figure}[htbp]
  \centering
  \includegraphics[width=0.90\textwidth]{figure_10_3_performance_map.png}
  \caption{Pace-reliability map for primary repeated batches with at least ten attempts. Vertical intervals are 95\% Wilson intervals for completion; horizontal positions are medians of completed laps. The map is a compact decision view rather than a claim of a controlled ablation between Agent~7 and Agent~8. Source: project evaluation records.}
  \label{fig:performance-map}
\end{figure}

Agent~8 produced the strongest observed learned result: 17/20 and 19/20 completions across two repeated batches, or 36/40 overall. Its best and median completed-lap time were both 83.038 seconds, indicating tightly clustered successful records but not erasing its four failures.

Agent~7 uses racing-line context, whereas Agent~8 uses sensors and history without an explicit route. Agent~8 therefore shows that a sensor-only policy can find an effective local strategy on \texttt{g-track-3}; it does not prove that racing-line information is unhelpful, because reward and training history also differ.

\subsection{Agent 8 version history}

\Cref{fig:agent8-development} retains the material design milestones behind the final result. It is a categorical evidence ladder, not a continuous training curve, because each point is a frozen policy after a different design decision. The v1 base completed 0/10; first completion reached 5/20 at 141.264 seconds; a pace candidate reached 2/20 at 87.299 seconds; and a stability candidate reached 5/20 at 83.118 seconds. Robust-pace and steering-limit trials both returned 0/20. The early self-imitation candidate completed 1/20 before the final champion reached 36/40 at 83.038 seconds.

\begin{figure}[htbp]
  \centering
  \includegraphics[width=\textwidth]{figure_10_4_agent8_development.png}
  \caption{Agent~8 material design milestones. Each categorical point is a separately evaluated frozen checkpoint or candidate; vertical intervals are 95\% Wilson intervals for completion. Crosses mark evaluations with no completed lap. The figure records changing designs and evidence, rather than implying a continuous common-budget learning curve. Source: project evaluation records.}
  \label{fig:agent8-development}
\end{figure}

The history shows that pace appeared before reliability and that apparently sensible interventions could fail under repeated evaluation. The final self-imitation and stability-oriented configuration was retained because it combined the fastest observed laps with materially stronger completion evidence. This remains a project-specific result on one track.

\subsection{Representative control behaviour}

Aggregate outcomes cannot explain how the controllers drove. \Cref{fig:control-profiles} therefore plots speed, lateral position and actions against normalised progress for three representative completed laps: the map-aware median (92.056 seconds), an Agent~7 v3 lap close to its 104.220-second median (104.180 seconds), and an 83.038-second Agent~8 lap. The traces are descriptive, resampled to common progress and lightly smoothed for legibility; raw telemetry remains stored.

\begin{figure}[htbp]
  \centering
  \includegraphics[width=\textwidth]{figure_10_5_control_profiles.png}
  \caption{Control profiles for representative completed laps. Speed, lateral position, steering and longitudinal action are aligned by normalised lap progress. The dotted lateral boundaries denote the track edge. Each line is a descriptive stored trace chosen near the relevant group median, resampled to a common grid and smoothed with a nine-point display window. Source: project telemetry records.}
  \label{fig:control-profiles}
\end{figure}

Agent~8 reaches higher speeds on several sections while remaining within the lateral boundaries in this trace. Agent~7 has larger steering reversals, consistent with its slower pace; the map-aware controller has a distinct longitudinal pattern because route and curvature are explicit. These profiles explain the stored outcomes but do not establish a causal ranking, since the policies differ in observation, reward and training history.

\subsection{Application result}

The evidence is not confined to scripts and model folders. The packaged application lets a user select an agent, start TORCS, inspect live telemetry, review a stored run and compare agents in one interface. Failure is informative rather than hidden: termination labels and telemetry distinguish stalling, reverse travel and leaving the circuit.

This supports the novice-facing objective: an AI newcomer can see why one fast lap is insufficient and compare rules, track information and learned policies without reading code. It is functional evidence of accessible deployment and explanation, not a claim of proven educational improvement.

